\documentclass[conference]{IEEEtran}
\IEEEoverridecommandlockouts
% The preceding line is only needed to identify funding in the first footnote. If that is unneeded, please comment it out.
\usepackage{cite}
\usepackage{amsmath,amssymb,amsfonts}
\usepackage{algorithmic}
\usepackage{graphicx}
\usepackage{textcomp}
\usepackage{multirow}
\usepackage{caption}
\usepackage{subcaption}
\usepackage{comment}
\usepackage{mwe}
\usepackage{xcolor}
\def\BibTeX{{\rm B\kern-.05em{\sc i\kern-.025em b}\kern-.08em
    T\kern-.1667em\lower.7ex\hbox{E}\kern-.125emX}}
\begin{document}

\title{Hybrid Misbehavior Detection for V2X: Compact Grid Features and Reward-Driven Q-Learning Against False Position Reports\\

}

\author{\IEEEauthorblockN{Chamath Gunawardena\IEEEauthorrefmark{1},
Owana Marzia Moushi\IEEEauthorrefmark{2},
Shengjie Xu\IEEEauthorrefmark{3}, and
Yi Qian\IEEEauthorrefmark{4},~\IEEEmembership}
\IEEEauthorblockA{\IEEEauthorrefmark{1}Department of Computer Science, Florida Polytechnic University, Lakeland, FL, USA}
\IEEEauthorblockA{\IEEEauthorrefmark{2}Department of Computer Science, University of Nebraska Omaha, Omaha, NE, USA}
\IEEEauthorblockA{\IEEEauthorrefmark{3}College of Information Science, University of Arizona, Tucson, AZ, USA}
\IEEEauthorblockA{\IEEEauthorrefmark{4}Tandy School of Computer Science, University of Tulsa, Tulsa, OK, USA}}

\maketitle

\begin{abstract}
Vehicular communication networks enable vehicles to exchange information through wireless links, enhancing road safety and supporting Intelligent Transportation Systems (ITS). Despite the use of cryptographic techniques to prevent external attacks, these networks remain vulnerable to insider threats where authenticated vehicles behave maliciously. One such threat is position falsification, in which attackers transmit false location information through Basic Safety Messages (BSMs), potentially causing unsafe driving decisions and traffic disruptions. Existing misbehavior detection systems often rely on high-dimensional feature sets, resulting in increased computational complexity and limited scalability. In this paper, we extend prior grid-based misbehavior detection frameworks by incorporating both supervised learning and reinforcement learning techniques for efficient attack detection. Our approach utilizes a lightweight grid-based representation of position information extracted from BSMs to reduce feature complexity while preserving discriminative capability. We first evaluate multiple supervised learning classifiers for multi-class detection of position falsification attacks, achieving high detection accuracy across diverse attack strategies. We then propose a reinforcement learning-based binary detection method using Q-learning for random and random-offset attacks, where each observer vehicle maintains an intelligent agent for neighboring transmitters and learns attack detection policies through reward-driven interactions with its environment. Rather than requiring fully pre-labeled training data, the RL framework uses unsupervised behavioral clustering and threshold-based reward validation to generate feedback during online learning. Experimental results on a publicly available dataset demonstrate that the supervised approach achieves strong detection performance, while the binary RL method provides improved adaptability for real-time deployment. The proposed methods outperform existing techniques while significantly reducing feature complexity.
\end{abstract}

\begin{IEEEkeywords}
Vehicular networks, grid-based misbehavior detection, network security, machine learning, reinforcement learning
\end{IEEEkeywords}

\section{Introduction}

Vehicular communication networks consist of wireless multi-hop connections with rapidly changing topology due to the high mobility of vehicles. These networks are a key component of Intelligent Transportation Systems (ITS), supporting traffic monitoring, transportation optimization, and accident reduction through vehicle-to-everything (V2X) communications~\cite{I1}. In vehicular networks, vehicles periodically transmit Basic Safety Messages (BSMs) containing information such as location, speed, acceleration, and heading. While cryptographic techniques can mitigate many external attacks, they remain ineffective against internal attacks involving authenticated malicious vehicles~\cite{I2}. As a result, misbehavior detection systems (MDSs) are widely used to identify malicious participants.

One critical internal threat is the position falsification attack, where malicious vehicles transmit incorrect position information through BSMs to mislead nearby vehicles or infrastructure. Common attack strategies include constant attacks, constant offset attacks, random attacks, random offset attacks, and eventual stop attacks.

Machine learning approaches have increasingly been adopted for detecting such attacks, and several ML-based MDSs have been proposed for position falsification detection. However, many existing methods rely on high-dimensional feature sets derived from raw BSM data, leading to increased computational complexity and memory overhead during training and deployment.


In this paper, we propose a reinforcement learning-based approach for detecting position falsification attacks in vehicular communication networks. To enable dynamic and online detection, we introduce a Q-learning-based framework in which observer vehicles maintain intelligent RL agents for neighboring transmitting vehicles within their communication range. These agents learn to identify malicious behavior through continuous interaction with the surrounding vehicular environment by analyzing \textcolor{red}{received BSMs} from monitored neighboring vehicles. Each RL agent maintains its own Q-table, which maps observed states, such as relative positions, velocities, and neighbor consistency metrics, to action values representing attack detection decisions.

\textcolor{red}{As new messages are received,} the Q-learning algorithm allows each RL agent to iteratively update its Q-values using reward signals derived from behavioral feedback indicating whether the monitored vehicle remains within a stable cluster associated with benign or malicious behavior, rather than relying on direct label supervision. Positive rewards are assigned when the selected action maintains cluster-consistent behavior, while negative rewards are assigned when the action causes unstable or inconsistent transitions, enabling the agent to gradually learn optimal attack detection policies through trial-and-error interactions. This reward-driven learning mechanism allows observer vehicles to continuously refine their detection strategies and adapt to dynamic traffic conditions or changing attacker behaviors over time. Unlike supervised learning methods, the reinforcement learning approach does not require a pre-collected labeled dataset or offline retraining, making it highly suitable for real-time deployment in practical vehicular communication environments. Furthermore, the current implementation of the proposed reinforcement learning algorithm effectively detects random attacks while providing the flexibility to be deployed in newly established networks without prior training data.


The supervised learning-based detection method \cite{Chamath1} leverages graph-based features extracted from a training dataset comprising position information collected from all transmitted \textcolor{red}{BSMs}. These features encode spatial and temporal relationships between vehicles, capturing topological and motion-based cues useful for distinguishing between legitimate and malicious behaviors. Although this approach benefits from detecting all five previously mentioned attack types, it requires access to the complete set of transmitted \textcolor{red}{messages} during the simulation period, which may not be feasible in real-time or online detection scenarios. \textcolor{red}{Consequently, while the supervised framework provides excellent offline detection performance using historical data, it is less suited for adaptive online deployment. This limitation motivates the reinforcement learning framework proposed in this paper, which continuously updates its detection policy through interaction with the vehicular environment and therefore complements the supervised approach by improving online adaptability.}

By integrating both detection methods, this study aims to evaluate the trade-offs between offline accuracy and online adaptability in identifying various forms of position falsification attacks in vehicular communication network environments.

\section{Related Work}

%In\cite{LR4}, the authors proposed an MDS using plausibility checks as a feature vector for machine learning to detect and classify misbehavior. The feature vector is 6 dimensional and includes 2 categorical features which are location plausibility checks and quantitative information on the vehicle's behavior.

Security in vehicular communication networks, particularly against position falsification and insider attacks, has been extensively studied in recent years. Since cryptographic methods cannot detect malicious behavior from authenticated vehicles, misbehavior detection systems (MDSs) have emerged as an essential complementary solution.

Early MDS approaches relied on rule-based and consistency-based techniques. Leinm\"uller et al.~\cite{leinmuller2006misbehavior} proposed plausibility checks using vehicle dynamics and communication constraints, while Raya et al.~\cite{raya2007data} introduced data-centric trust mechanisms based on message cross-verification. Although effective in simple scenarios, these methods struggle in dynamic environments due to their dependence on predefined rules.

More recently, machine learning-based approaches have gained attention. Grover et al.~\cite{grover2018mdp} modeled misbehavior detection as a sequential decision-making problem using Markov Decision Processes, demonstrating the potential of adaptive learning in vehicular environments. Supervised learning methods have also been widely explored, with Van der Heijden et al.~\cite{heijden2018survey} highlighting the effectiveness of high-dimensional feature extraction from \textcolor{red}{BSMs}. However, such methods often require extensive feature engineering and large training datasets, limiting real-time applicability.

To improve efficiency, lightweight and spatial consistency-based methods have also been proposed~\cite{bissmeyer2011trustworthy,sedjelmaci2018intrusion}, though balancing detection accuracy, computational complexity, and adaptability remains challenging. In this work, we build upon these efforts by investigating reinforcement learning and supervised learning approaches for position falsification detection using reduced-complexity position-based features.

\section{System Model, Dataset, and Attack Models}

\subsection{System Model}\label{AA}

There are three types of MDS in vehicular networks~\cite{I3}: standalone, cooperative/distributed, and hierarchical. In a standalone MDS, each vehicle independently detects misbehavior using only its locally collected data without cooperation from other vehicles. In a cooperative/distributed MDS, vehicles collaborate by sharing detection results, and decisions are made through distributed aggregation. In a hierarchical MDS, vehicles are organized into clusters, where each cluster elects a cluster head that aggregates detection results from cluster members to make final decisions.

Specifically, each vehicle maintains local RL agents that observe temporal BSM position sequences received from neighboring transmitters under evaluation, with one RL agent instantiated per observed transmitter. The local decision of each agent is forwarded to the cluster head, which aggregates per-vehicle decisions from surrounding observers to produce the final cooperative detection output.

In our system model, we propose a cooperative grid-based machine learning scheme for detecting position falsification attacks. Each vehicle is equipped with an MDS that makes detection decisions based on BSM data collected from nearby vehicles, and the aggregated results are used to identify and remove misbehaving vehicles from the network.

\begin{figure}[htbp]
\centerline{\includegraphics[scale=.6]{systemarchitecture.png}}
\caption{Vehicular Network Model for ML-based MDS}
\label{fig6_0}
\end{figure}

As shown in Figure~\ref{fig6_0}, the system architecture consists of a Trusted Authority (TA), Access Points (APs), vehicles, and a misbehavior detection system (MDS). The TA acts as a central trusted entity responsible for vehicle registration, key management, and verification. APs facilitate communication between the TA and vehicles via cellular infrastructure. Each vehicle is equipped with an MDS that actively detects malicious behavior upon receiving BSMs. Vehicles periodically broadcast BSMs containing information such as location, time, and speed, and report detection results to the TA through APs.

\subsection{Dataset}

The VeReMi dataset~\cite{veremi} is a publicly available benchmark for evaluating misbehavior detection in vehicular networks. It is generated using LuST~\cite{lust} and Veins~\cite{veins}. LuST models realistic traffic in Luxembourg using the Simulation of Urban Mobility (SUMO)~\cite{sumo}, while Veins integrates SUMO with OMNeT++~\cite{omnet} for network simulation.

The dataset contains message logs from multiple simulations, including GPS data and received BSMs. It covers five position falsification attacks, three vehicle densities (Low, Medium, High), and three attacker ratios ($10\%$, $20\%$, $30\%$), with each configuration repeated five times, resulting in 225 simulations.

Each \textcolor{red}{message entry} includes timestamps, sender information, message ID, position and speed vectors, RSSI, and noise vectors. Corresponding ground truth files provide actual vehicle information, including true position, speed, and attacker type, enabling comparison with potentially falsified data in the logs.

\subsection{Attacker Models}

The VeReMi dataset includes five types of position falsification attacks: constant, constant offset, random, random offset, and eventual stop attacks. In our grid-based model, vehicle positions from transmitted BSMs are mapped onto a grid to identify anomalous patterns. 



\begin{figure*}[htbp]
\centering
\begin{subfigure}{0.3\textwidth}
\centering
\includegraphics[width=\linewidth]{Figure_1N.png}
\caption{Normal trajectory}
\end{subfigure}
\begin{subfigure}{0.3\textwidth}
\centering
\includegraphics[width=\linewidth]{Figure_A4.png}
\caption{Random attack}
\end{subfigure}
\begin{subfigure}{0.3\textwidth}
\centering
\includegraphics[width=\linewidth]{Figure_A16.png}
\caption{Eventual stop attack}
\end{subfigure}
\caption{Representative spatial transmission patterns used to motivate the proposed detection features. Normal trajectories exhibit continuous motion patterns, random attacks exhibit spatial inconsistency, and eventual-stop attacks show repeated stationary transmissions after initial motion. These behavioral differences motivate the grid occupancy and repetition-based detection features used in the proposed framework.}
\label{fig_attack_compact}
\end{figure*}

\begin{table}[htbp]
\caption{Summary of Position Falsification Attack Types}
\centering
\renewcommand{\arraystretch}{1.2}
\begin{tabular}{|p{2.1cm}|p{5.8cm}|}
\hline
\textbf{Attack Type} & \textbf{Behavioral Description} \\ \hline
Constant / Constant Offset & Attacker repeatedly transmits either a fixed false position or a consistently shifted trajectory, creating low spatial variance but abnormal positional consistency. \\ \hline
Random / Random Offset & Attacker transmits arbitrary false positions or randomly perturbed offsets, producing highly inconsistent spatial trajectories. \\ \hline
Eventual Stop & Attacker initially behaves normally before repeatedly transmitting the same final position to mimic a stopped vehicle. \\ \hline
\end{tabular}
\label{tab:attack_summary}
\end{table}

The attack models considered in this work include several common forms of position falsification. A \textit{constant attack} occurs when a vehicle repeatedly transmits the same fixed position regardless of its true movement. In a \textit{constant offset attack}, the vehicle reports its actual trajectory with a consistent positional offset applied. A \textit{random attack} involves broadcasting arbitrary positions with no spatial consistency, while a \textit{random offset attack} applies varying random offsets to the true position within a constrained range. Lastly, an \textit{eventual stop attack} occurs when a vehicle initially behaves normally before repeatedly transmitting its final valid position, falsely indicating that it has stopped moving.

Table~\ref{tab:attack_summary} summarizes the representative behavioral characteristics of the considered attack classes and highlights why these patterns are distinguishable using the proposed grid-based spatial features.



\section{Reinforcement Learning-Based Misbehavior Detection in Vehicular Communication Networks}


In this section, we propose a reinforcement learning (RL)-based scheme to detect position falsification attacks in vehicular networks. Our approach enables each observer vehicle to maintain RL agents for neighboring transmitting vehicles within its communication range. Each RL agent monitors and analyzes the \textcolor{red}{received BSMs} associated with a neighboring vehicle to assess the likelihood of malicious behavior, with particular focus on identifying falsified position reports.


\subsection{System Overview}


Each observer vehicle $v_o$ in the VANET maintains a set of RL agents, where each agent $A_i$ is associated with a neighboring transmitting vehicle $v_i$ being monitored. Rather than monitoring its own transmissions, the observer vehicle assigns an RL agent to each neighboring transmitter from which BSMs are received. The RL agent stores the history of \textcolor{red}{received BSMs} transmitted by the monitored neighboring vehicle, with each \textcolor{red}{message} containing positional data (e.g., GPS coordinates). Let $P_i = \{p_{i}^{(1)}, p_{i}^{(2)}, \dots, p_{i}^{(n_i)}\}$ denote the sequence of positions, where $p_{i}^{(j)} = (x_{i}^{(j)}, y_{i}^{(j)})$, reported in \textcolor{red}{messages} received from neighboring vehicle $v_i$ over time, and where $n_i$ is the number of \textcolor{red}{messages} observed from $v_i$ by the observer vehicle $v_o$.


\subsection{Average Distance Calculation}

For each vehicle $v_i$, if $n_i < 2$, the average distance is set to zero:
\[
D_i = 0 \quad \text{if } n_i < 2
\]

Otherwise, the average distance $D_i$ traveled by $v_i$ is calculated using:
\[
D_i = \frac{1}{n_i - 1} \sum_{j=1}^{n_i - 1} \|p_{i}^{(j+1)} - p_{i}^{(j)}\|_2
\]
where $\| \cdot \|_2$ denotes the Euclidean distance between two consecutive position points.

\subsection{Clustering Based on Average Distance}


At each RL training round, clustering is performed on the set of vehicles with $n_i \geq 2$ using their computed average distances $\{D_i\}$ as the input feature to partition observations into two candidate groups. Specifically, K-means clustering with $k=2$ is applied to separate vehicles based on their average inter-report distances. After clustering, the average inter-report distance threshold $\phi$ is applied to interpret and validate the resulting clusters as either $\mathcal{C}_{\text{low}}$ or $\mathcal{C}_{\text{high}}$. The two resulting clusters are defined as:
\begin{itemize} 
    \item $\mathcal{C}_{\text{low}}$: Vehicles with relatively lower average distances.
    \item $\mathcal{C}_{\text{high}}$: Vehicles with relatively higher average distances.
\end{itemize}


\subsection{Reinforcement Learning Formulation}

Each agent $A_i$ uses a Q-learning algorithm, where the RL environment is defined as follows:

\begin{itemize}
    \item \textbf{State space} $\mathcal{S}$: 
    \[
    \mathcal{S} = \{s_{\text{low}}, s_{\text{high}}\}
    \]
    where $s_{\text{low}}$ indicates the vehicle belongs to $\mathcal{C}_{\text{low}}$, and $s_{\text{high}}$ indicates it belongs to $\mathcal{C}_{\text{high}}$.

    \item \textbf{Action space} $\mathcal{A}$:
    \[
    \mathcal{A} = \{\text{stay}, \text{switch}\}
    \]
    where \texttt{stay} indicates that the vehicle remains in the same cluster as the previous round, and \texttt{switch} indicates the observed cluster assignment has changed between consecutive rounds.

    \item \textbf{Reward function} $R(s,a)$:
    The reward function is designed to reflect suspicion based on vehicle behavior:
    \[
    R(s, a) = 
    \begin{cases}
        +1, & \text{if } s = s_{\text{low}} \text{ and } a = \text{stay} \\
        0,  & \text{if } s = s_{\text{high}} \text{ and } a = \text{stay} \\
        -1, & \text{if } a = \text{switch} \\
    \end{cases}
    \]
    This reward structure encourages stability in the low-distance cluster while penalizing frequent switching or persistent presence in the high-distance cluster, which are indicative of suspicious behavior. In the implementation, the transition from one clustering round to the next is treated as an observed stay/switch outcome, and the reward is assigned according to the paper's mapping: $+1$ for low/stay, $0$ for high/stay, and $-1$ otherwise.
\end{itemize}

\subsection{Q-Learning Algorithm}

The Q-value update rule for each agent is given by:

\begin{equation*}
\begin{aligned}
Q(s_t, a_t) \leftarrow Q(s_t, a_t) + \alpha \Big[ R(s_t, a_t) \\
+ \gamma \max_{a'} Q(s_{t+1}, a') - Q(s_t, a_t) \Big]
\end{aligned}
\end{equation*}

where:
\begin{itemize}
    \item $s_t$ and $a_t$ are the current state and action.
    \item $\alpha \in (0, 1]$ is the learning rate.
    \item $\gamma \in [0, 1]$ is the discount factor.
    \item $Q(s, a)$ is the Q-value representing the expected cumulative reward for taking action $a$ in state $s$.
\end{itemize}

\subsection{Misbehavior Classification}

After multiple training rounds, each agent maintains a Q-value profile for its vehicle. We derive a behavior score for each vehicle based on the cumulative Q-values:
\[
\text{Score}(v_i) = Q(s_{\text{low}}, \text{stay}) - Q(s_{\text{high}}, \text{stay}) - Q(s_{\text{high}}, \text{switch})
\]

This score captures the agent’s preference for the low-distance cluster and penalizes actions associated with suspicious behavior. A threshold $\theta$ is defined such that:
\[
\text{Vehicle } v_i =
\begin{cases}
\text{benign}, & \text{if } \text{Score}(v_i) > \theta \\
\text{malicious}, & \text{otherwise}
\end{cases}
\]

The threshold $\theta$ can be determined empirically or via a validation dataset of known benign/malicious behaviors.

This RL-based detection mechanism provides a decentralized and adaptive method for detecting vehicles that engage in position falsification. By leveraging average motion characteristics and reinforcement learning, each vehicle's behavior is evaluated locally, reducing reliance on centralized monitoring while improving resilience against sophisticated adversaries.

% EMPHASIZE THAT WE ARE DETECTING RANDOM OR RANDOM OFFSET ATTACKS. BINARY CLASSIFICATION RANDOM ATTACK VS BENIGN
The reinforcement learning-based detection method is specifically applied to detect \textit{random} or \textit{random offset} position falsification attacks in vehicular networks. In these types of attacks, malicious vehicles either transmit completely arbitrary position values (random attack) or report locations offset by randomly varying vectors from their true positions (random offset attack). These behaviors can cause significant disruption in position-based vehicular applications and require dynamic detection mechanisms.

This online and distributed detection mechanism enables real-time identification of vehicles exhibiting anomalous behaviors consistent with random or random offset attacks, without requiring full knowledge of the global traffic state or a labeled training dataset. The binary decision output contributes to the overall system's resilience by enabling timely isolation of suspected malicious nodes in the network.

\begin{figure}[htbp]
\centerline{\includegraphics[width=1\linewidth]{MaliciousC.png}}
\caption{Q-value distribution grid for a client suspected of position falsification attack}
\label{qgrid}
\end{figure}



\section{Graph based Supervised Learning}

% MENTION THAT RESULTS CAN BE FOUND IN CONFERENCE PAPER REFERENCE
In our previous scheme\cite{Chamath1}, we transform the data by taking the positional data from all the BSMs such that each sample represents a vehicle instead to capture the vehicle's positional pattern on the road. Any unusual positions of vehicles where they seem to be well off road would be captured and marked as misbehaving. The positional data for each vehicle is represented on our grid-based scheme in which we create seven features. The proposed MDS involves a grid-based approach to generate seven new features along with one other feature taken from the BSMs. The grid-based MDS uses these features on a multi-class classifier to classify and detect which vehicles are committing the position falsification attacks.

A rectangular grid is created for each vehicle by plotting all of the transmitted positions from the BSMs for each vehicle. A 2-D matrix is created to represent a vehicle's grid and capture the vehicle's transmitted positions. 

\begin{figure}[htbp]
\centerline{\includegraphics[width=1\linewidth]{gridEX.PNG}}
\caption{10 x 10 Grid of vehicle committing eventual stop attack}
\label{grid}
\end{figure}

To normalize our features, the grid for each vehicle needs to be in the same scale on both axes. To address this, let $x_{min}$ be the minimum x-coordinate out of all the transmitted positions in the dataset, let $x_{max}$ be the maximum x-coordinate out of all the transmitted positions in the dataset, let $y_{min}$ be the minimum y-coordinate out of all the transmitted positions in the dataset and let $y_{max}$ be the maximum y-coordinate out of all the transmitted positions in the dataset. The coordinates: $\{$($x_{min}, y_{min}$), ($x_{min}, y_{max}$), ($x_{max}, y_{min}$), ($x_{max}, y_{max}$) $\}$ are used as the corners of the rectangular grid. The size of each window by the x-axis is represented by $\Delta x$ and the size of each window by the y-axis is represented by $\Delta y$. Therefore, the size of each window in a grid is $\Delta x$ by $\Delta y$. To create an $n \times n$ grid, we calculate $\Delta x$ and $\Delta y$ as follows:

\begin{equation}
    \Delta x = \frac{x_{max} - x_{min}}{n}
\end{equation}

\begin{equation}
    \Delta y = \frac{y_{max} - y_{min}}{n}
\end{equation}

Let $K$ represent the set of all vehicles in the dataset, and $V^k$ ($k \in K$) represent an $n \times n$ matrix in which, the $(i,j)$th entry of $V^k$, $v^k_{i,j}$, represents the total number of transmitted positions for the window in the $i$th row and $j$th column of the grid.

Let $M^k$ be the set of all messages transmitted by the vehicle $k \in K$. We initialize $V^k$ to a zero matrix. For each message $m \in M^k$, let $(x_m,y_m)$ be the BSM transmitted position. We can find the corresponding window by calculating $(i,j)$ as follows: 
\begin{equation}
    i = \left \lfloor{\frac{y_m - y_{min}}{\Delta y}}\right \rfloor 
\end{equation}
\begin{equation}
    j = \left \lfloor{\frac{x_m - x_{min}}{\Delta x}}\right \rfloor 
\end{equation}

We increment $v^k_{i,j}$ by $1$.

Figure~\ref{grid} shows a $10$ x $10$ matrix corresponding to the grid of a vehicle's positions in a simulation. The grid shows that the first three BSMs from this vehicle are legitimate with three different positions transmitted once each but the fourth position is transmitted $89$ times representing an eventual stop attack.

\subsection{Proposed Features}

Rather than using the message log data, we use the grid data for each vehicle and propose new features to greatly reduce the size of the dataset. 

The first new feature, $p_k$, corresponds to the total number of positions from all BSMs transmitted for any vehicle $k$ (i.e., the sum of all the entries in $V^k$):

\begin{equation}
    p_k = \sum_{i = 0}^{n-1} \sum_{j = 0}^{n-1} v^k_{i,j}
\end{equation}

The second new feature, $w_k$, represents the number of windows in any vehicle $k$'s grid that contains points. Let $a$ be an integer. We define a function $X$ such that $X(a) = 1$, if $a > 0$, and $X(a) = 0$ otherwise.

\begin{equation}
    w_k = \sum_{i = 0}^{n-1} \sum_{j = 0}^{n-1} X(v^k_{i,j})
\end{equation}

The third feature is called the spread ratio ($sr$) which is between the second (number of windows containing points) and first feature (total points) for any vehicle $k$ ($sr_k = \frac{w_k}{p_k}$). $sr$ of a vehicle captures how spread out the vehicle's transmitted positions are on its grid. 

We define a binary $n \times n$ attacking matrix $A$ for the grid using the whole training set. We set $a_{i,j} = 0$ if $v^k_{i,j} > 0 \hspace{0.9mm} \forall k \in K$ where $v^k_{i,j}$ has at least one position from a benign vehicle, and $a_{i,j} = 1$ otherwise. The windows with $a_{i,j} = 1$ are considered attacking windows which may include empty windows. The fourth feature, $q$ corresponds to the total number of points in attacking windows for any vehicle $k$'s grid.

\begin{equation}
    q = \sum_{i = 0}^{n-1} \sum_{j = 0}^{n-1} v^k_{i,j} a_{i,j} 
\end{equation}

The fifth feature is called the attack ratio ($ar_k$) which is the ratio between the fourth (total number of points in attacking windows) and first feature (total points) for any vehicle $k$ ($ar_k = \frac{q_k}{p_k}$). With the attack ratio, the vehicle positions outside of the windows containing positions from benign vehicles are captured such as positions that are clearly not on the road.  

The sixth feature represents the average absolute difference between all consecutive points for any vehicle $k$.

Finally, the seventh feature corresponds to the average absolute speed from all the BSMs sent from any vehicle $k$. Any parked vehicle will be differentiated from a constant attack with this feature because its average speed will be $0$, while a constant attacking vehicle will have its actual speed transmitted.

 

\section{Performance Evaluation}

To evaluate the proposed reinforcement learning MDS system, we dynamically store data from transmitted \textcolor{red}{messages}, sequentially, via timestamp and run the detection scheme. The model is trained for binary classification for the random and random offset attacks from the VeReMi dataset. 

To evaluate the proposed grid-based MDS system, we have trained and tested on five classification models using the VeReMi dataset. The models were trained and tested with multi-classification and compared to other existing approaches that detected position falsification attacks on the VeReMi dataset using multi-classification.

\subsection{Evaluation Metrics}

The two main metrics used to evaluate the performance of the proposed scheme are model accuracy and F1 score. The model accuracy is the ratio between the number of correctly predicted observations and the number of total observations.

\begin{equation}
    Accuracy = \frac{TP + TN}{TP + TN + FP + FN}
\end{equation}

Where $TP$ is the number of true positives, $TN$ is the number of true negatives, $FP$ is the number of false positives and $FN$ is the number of false negatives.

Precision is the ratio between the number of correctly predicted malicious vehicles and the total number of predicted malicious vehicles: $\frac{TP}{TP + FP}$. Recall is the ratio between the number of correctly predicted malicious vehicles and the total number of actual malicious vehicles: $\frac{TP}{TP + FN}$. The F1 score is the harmonic mean of precision and recall. 

\begin{equation}
    F1 score = 2 \times \frac{precision \times recall}{precision + recall}
\end{equation}

\begin{figure}[htbp]
\centerline{\includegraphics[width=1\linewidth]{SmallR.png}}
\caption{Small Simulation Random Attack Detection}
\label{smallr}
\end{figure}

\begin{figure}[htbp]
\centerline{\includegraphics[width=1\linewidth]{MediumR.png}}
\caption{Medium Simulation Random Attack Detection}
\label{mediumr}
\end{figure}

\begin{figure}[htbp]
\centerline{\includegraphics[width=1\linewidth]{SmallRO.png}}
\caption{Small Simulation Random Offset Attack Detection}
\label{smallro}
\end{figure}

\begin{figure}[htbp]
\centerline{\includegraphics[width=1\linewidth]{MediumRO.png}}
\caption{Medium Simulation Random Offset Attack Detection}
\label{mediumro}
\end{figure}

To assess the effectiveness of the proposed reinforcement learning-based detection scheme for position falsification attacks in vehicular networks, we conducted experiments using the VeReMi dataset. The detection performance was evaluated against two types of attacks: \textit{random} and \textit{random offset}, across two simulation scales: \textit{small} and \textit{medium}. Detection accuracy was plotted against the number of \textcolor{red}{received messages}, with four distinct experimental conditions captured in Figures~\ref{smallr}, \ref{mediumr}, \ref{smallro}, and \ref{mediumro}.

\subsection{Detection Performance in Small-Scale Simulations}

In the small-scale simulation environment, the system achieved high detection performance with relatively few messages:

\begin{itemize}
    \item \textbf{Random Attack:} As shown in Figure~\ref{smallr}, the detection accuracy exceeded 90\% after only \textcolor{red}{20 received messages}. Beyond this point, the accuracy continued to increase and remained consistently above 90\%, demonstrating rapid and stable learning behavior.
    
    \item \textbf{Random Offset Attack:} In contrast, Figure~\ref{smallro} shows that detecting the more subtle random offset attack required approximately \textcolor{red}{400 received messages} to surpass the 90\% detection threshold. After this point, accuracy also continued to improve and remained steadily high.
\end{itemize}

\subsection{Detection Performance in Medium-Scale Simulations}

The proposed detection scheme maintained strong performance under medium-scale simulation conditions:

\begin{itemize}
    \item \textbf{Random Attack:} As depicted in Figure~\ref{mediumr}, the detection accuracy crossed the 90\% mark after approximately \textcolor{red}{100 received messages}. The model maintained and improved upon this accuracy throughout the remaining transmissions.
    
    \item \textbf{Random Offset Attack:} For this more challenging scenario, Figure~\ref{mediumro} indicates that the 90\% threshold was reached after about \textcolor{red}{500 received messages}. Detection performance remained stable and increased slightly thereafter.
\end{itemize}


The results demonstrate that the reinforcement learning-based detection agent is highly effective in identifying both random and random offset position falsification attacks in vehicular networks. The detection scheme responds quickly to obvious anomalies (e.g., random attacks), achieving high accuracy with minimal BSMs, while still effectively adapting to subtler adversarial behaviors like random offset attacks with slightly longer observation periods. These findings validate the scalability and adaptability of our approach across different traffic densities and attack types.



The reinforcement learning hyperparameters used in our experiments are summarized as follows: learning rate $\alpha = 0.1$, discount factor $\gamma = 0.9$, exploration rate $\epsilon = 0.1$, and decision threshold $\theta = 0.5$. The Q-learning agent was trained for 500 episodes using \textcolor{red}{a trace-level chronological 70/30 split of the VeReMi message stream}, where \textcolor{red}{the first 70\% of the trace files are used for training and the remaining 30\% for evaluation}. Clustering was used to partition observations into two candidate groups, and the average inter-report distance threshold $\phi$ is data-dependent and selected during validation to interpret and validate the resulting clusters as either $C_{low}$ or $C_{high}$. In terms of computational overhead, the Q-table requires $|S|\times|A|$ entries where $|S|$ is the number of discretized states and $|A|=2$ actions, resulting in lightweight memory requirements suitable for onboard vehicular deployment. \textcolor{red}{To suppress artificial oscillations caused by late-arriving senders and short-run turnover in the sampled stream, the reported RL accuracy curve is smoothed with a rolling window over the sampled accuracy points. This produces a gradual ascent with small fluctuations and a stable plateau rather than a sharp jump from 0\% to 100\%. In addition, the evaluation metric is computed in a cumulative, evidence-based manner with a minimum of five BSMs per sender before a decision is counted, so accuracy improves as more BSM evidence is accumulated and remains stable or slightly increasing at the end of the simulation rather than declining after a short burst of perfect predictions. For senders not yet represented in the learned Q-table, the policy defaults to the benign score until enough evidence is accumulated, preventing an unseen vehicle from being misclassified as malicious by default. Per-message processing} consists only of state lookup and Q-value update, resulting in constant-time complexity $O(1)$ after state construction.


\begin{table*}[]
\begin{center}

\begin{tabular}{ |c|c|c|c|c| } 
\hline
Method & Task Type & Attack Types & Reported Metric & Source\\
\hline
Ref~\cite{LLR1} & Multi-class & Multiple & F1-score: 0.744 & Prior Work\\ 
Ref~\cite{LLR7} & Multi-class & Multiple & F1-score: 0.954 & Prior Work\\ 
Ref~\cite{LLR10} & Multi-class & Multiple & F1-score: 0.985 & Prior Work\\ 
Ref~\cite{Chamath1} & Multi-class & Multiple & F1-score: 0.992 & Prior Work\\ 
Proposed Supervised Grid-Based & Multi-class & All VeReMi Attacks & F1-score: 0.996 & This Work\\ 
Proposed RL Q-Learning & Binary & Random / Random Offset & Accuracy: See Fig. \ref{smallr}--\ref{mediumro} & This Work\\ 
\hline
\end{tabular}

\caption{Comparison of supervised grid-based and reinforcement learning results with prior approaches. Reported metrics are presented as originally provided in the cited works.}
\label{tab1}
\end{center}
\end{table*}

Table~\ref{tab1} shows a comparison of the proposed scheme with other approaches in detecting position falsification attacks using the VeReMi dataset. The F1-score of 0.996 corresponds to the proposed supervised multi-class grid-based classifier evaluated using grid resolution $n=60$ (number of windows = 3600), which was found empirically to provide the best performance. In contrast, the reinforcement learning framework is evaluated separately in Figures~\ref{smallr}--\ref{mediumro} under a binary classification setting for random and random offset attacks and is therefore not directly merged into the scalar F1 comparison table.



\section{Conclusion}


We proposed a reinforcement learning-based detection framework for identifying position falsification attacks, specifically random and random offset attacks, in vehicular networks. By leveraging Q-learning, the proposed framework enables each observer vehicle to maintain intelligent RL agents for neighboring transmitters, where each agent continuously improves its attack detection policy through reward-driven interactions with the environment, allowing adaptive and autonomous identification of malicious behavior without requiring prior labeled training data. Using the VeReMi dataset, we evaluated the effectiveness of our approach across both small and medium-scale simulation environments.
The experimental results demonstrated that the proposed method achieves high detection accuracy with a relatively small number of \textcolor{red}{received messages}. For random attacks, the detection system reached over 90\% accuracy within the first \textcolor{red}{20 to 100 received messages}, depending on simulation scale. For the more subtle random offset attacks, the system required a larger number of messages, approximately \textcolor{red}{400 to 500 received messages}, to achieve similar accuracy levels. In all cases, once the 90\% threshold was surpassed, the detection accuracy remained consistently high and continued to improve as more messages were observed, indicating that the reinforcement learning agent progressively refines its detection strategy over time as additional environmental feedback is received.
These findings highlight the adaptability, robustness, and online learning capability of the reinforcement learning agent in dynamic vehicular communication network scenarios, particularly its ability to adjust to evolving attack patterns and varying network conditions. Future work will focus on extending the detection framework to more complex attack models, integrating additional environmental context, and evaluating performance under real-world driving conditions.


\section*{Acknowledgement}

This work was partially supported by National Science Foundation under grants CNS-2008145, CNS-2007995, CNS-2319486, CNS-2319487, CNS-2344341.

%\bibliographystyle{abbrv}
\bibliographystyle{unsrt}
\bibliography{References}

\end{document}
